\documentclass[a4paper,10pt,twocolumn]{ltjarticle}
\usepackage{kaitoushu}

% グラフ用の共通スタイル定義
\pgfplotsset{
  myaxis/.style={
    axis lines=middle,
    xlabel={$x$}, ylabel={$y$},
    every axis x label/.style={at={(ticklabel* cs:1)},anchor=west},
    every axis y label/.style={at={(ticklabel* cs:1)},anchor=south},
    tick label style={font=\scriptsize},
    label style={font=\small},
    width=5.5cm, height=5cm,
    grid=none,
    samples=100,
  }
}

\begin{document}

\problemrange{4章：指数関数と対数関数 — Basic (p.49) 対数（問225--239）}



\mondai{225}{次の値を求めよ．}
\kotae{
対数の定義：$\log_{a}b=c$ $\Leftrightarrow$ $a^{c}=b$．「$a$ を何乗すると $b$ になるか」．

(1) $\log_{5}25=\log_{5}5^{2}=2$

(2) $\log_{3}\dfrac{1}{27}=\log_{3}3^{-3}=-3$

(3) $\log_{0.1}0.001=\log_{0.1}(0.1)^{3}=3$

(4) $\log_{4}1=0$\quad（任意の底で $\log_{a}1=0$）

(5) $\log_{2}0.25=\log_{2}\dfrac{1}{4}=\log_{2}2^{-2}=-2$

(6) $\log_{2}\sqrt[5]{8}=\log_{2}(2^{3})^{1/5}=\log_{2}2^{3/5}=\dfrac{3}{5}$
}

\mondai{226}{次の式を計算せよ．}
\kotae{
対数の性質：$\log_{a}MN=\log_{a}M+\log_{a}N$，$\log_{a}\dfrac{M}{N}=\log_{a}M-\log_{a}N$，$\log_{a}M^{k}=k\log_{a}M$．

(1) $\log_{2}64=\log_{2}2^{6}=6$

(2) $\log_{3}6+\log_{3}\dfrac{3}{2}=\log_{3}\!\left(6\cdot\dfrac{3}{2}\right)=\log_{3}9=2$

(3) $\log_{5}2-\log_{5}10=\log_{5}\dfrac{2}{10}=\log_{5}\dfrac{1}{5}=\log_{5}5^{-1}=-1$

(4) $\dfrac{1}{2}\log_{2}10+\log_{2}\dfrac{4}{\sqrt{5}}$

$=\log_{2}\sqrt{10}+\log_{2}\dfrac{4}{\sqrt{5}}=\log_{2}\!\left(\sqrt{10}\cdot\dfrac{4}{\sqrt{5}}\right)=\log_{2}(4\sqrt{2})=\log_{2}2^{5/2}=\dfrac{5}{2}$
}

\mondai{227}{等式 $\log_{a}\dfrac{1}{\sqrt[n]{M}}=-\dfrac{1}{n}\log_{a}M$ を証明せよ．}
\kotae{
左辺$=\log_{a}M^{-1/n}=-\dfrac{1}{n}\log_{a}M=$右辺

（$\log_{a}M^{k}=k\log_{a}M$ を使った）\qed
}

\mondai{228}{$A,B,C,a$ は正の数で $a\neq1$ とするとき，次の問いに答えよ．}
\kotae{
(1) $\log_{a}\dfrac{A}{B}+\log_{a}\dfrac{B}{C}+\log_{a}\dfrac{C}{A}=0$ を証明．

左辺$=\log_{a}\!\left(\dfrac{A}{B}\cdot\dfrac{B}{C}\cdot\dfrac{C}{A}\right)=\log_{a}1=0=$右辺 \qed

(2) $\log_{3}\sqrt{6}+\log_{3}\sqrt{10}-\log_{3}\sqrt{20}$

$=\log_{3}\dfrac{\sqrt{6}\cdot\sqrt{10}}{\sqrt{20}}=\log_{3}\sqrt{\dfrac{60}{20}}=\log_{3}\sqrt{3}=\dfrac{1}{2}$
}

\mondai{229}{$\log_{a}2$ の値を求めよ．}
\kotae{
問題文から底 $a$ が指定されていない場合，底の変換公式等を利用する．

※ この問題は教科書 p.115 を参照．底が $10$ の場合：$\log_{10}2\approx0.3010$（常用対数）．
}

\mondai{230}{1でない正の数 $a,b,c$ について，次の式を計算せよ．$(\log_{a}b)\cdot(\log_{b}c)\cdot(\log_{c}a)$}
\kotae{
底の変換公式：$\log_{a}b=\dfrac{\log b}{\log a}$（底は何でもよい）

$(\log_{a}b)(\log_{b}c)(\log_{c}a)=\dfrac{\log b}{\log a}\cdot\dfrac{\log c}{\log b}\cdot\dfrac{\log a}{\log c}=1$

（分子と分母がすべて約分される）
}

\mondai{231}{次の式を簡単にせよ．}
\kotae{
底の変換公式で底をそろえてから計算する．

(1) $(\log_{2}27)\cdot(\log_{9}8)$

$=\log_{2}3^{3}\cdot\log_{9}2^{3}=3\log_{2}3\cdot3\log_{9}2=9\log_{2}3\cdot\dfrac{\log 2}{\log 9}$

$=9\cdot\dfrac{\log 3}{\log 2}\cdot\dfrac{\log 2}{2\log 3}=9\cdot\dfrac{1}{2}=\dfrac{9}{2}$

(2) $(\log_{3}125)\cdot(\log_{4}3)\cdot(\log_{5}32)$

$=\dfrac{\log 125}{\log 3}\cdot\dfrac{\log 3}{\log 4}\cdot\dfrac{\log 32}{\log 5}=\dfrac{3\log 5}{\log 3}\cdot\dfrac{\log 3}{2\log 2}\cdot\dfrac{5\log 2}{\log 5}=\dfrac{3\cdot5}{2}=\dfrac{15}{2}$
}

\mondai{232}{次の関数のグラフをかけ．}
\kotae{
対数関数 $y=\log_{a}x$ の基本性質：必ず点$(1,0)$を通り，$a>1$ なら右上がり，$0<a<1$ なら右下がり．

(1) $y=\log_{3}x$

\begin{tikzpicture}
\begin{axis}[myaxis, xmin=-1,xmax=10,ymin=-3,ymax=3,
  xtick={1,3,9}]
\addplot[thick,blue,domain=0.05:9.5]{ln(x)/ln(3)};
\addplot[only marks,mark=*,mark size=1.5pt] coordinates {(1,0)(3,1)(9,2)};
\end{axis}
\end{tikzpicture}

(2) $y=\log_{1/3}x$

底 $\dfrac{1}{3}<1$ なので右下がり．$y=\log_{3}x$ の$x$軸対称．

\begin{tikzpicture}
\begin{axis}[myaxis, xmin=-1,xmax=10,ymin=-3,ymax=3,
  xtick={1,3,9}]
\addplot[thick,blue,domain=0.05:9.5]{ln(x)/ln(1/3)};
\addplot[only marks,mark=*,mark size=1.5pt] coordinates {(1,0)(3,-1)};
\end{axis}
\end{tikzpicture}

(3) $y=\log_{4}(x+1)$

$y=\log_{4}x$ を$x$方向に$-1$平行移動．点$(0,0)$を通る．

\begin{tikzpicture}
\begin{axis}[myaxis, xmin=-2,xmax=8,ymin=-2,ymax=2.5]
\addplot[thick,blue,domain=-0.95:7.5]{ln(x+1)/ln(4)};
\addplot[dashed,red] coordinates {(-1,-2)(-1,2.5)};
\addplot[only marks,mark=*,mark size=1.5pt] coordinates {(0,0)};
\node[font=\scriptsize] at (axis cs:-1.5,1.5) {$x=-1$};
\end{axis}
\end{tikzpicture}

(4) $y=\log_{4}(-x)$

$y=\log_{4}x$ を$y$軸対称（$x\to-x$）．$x<0$ で定義．

\begin{tikzpicture}
\begin{axis}[myaxis, xmin=-8,xmax=2,ymin=-2,ymax=2.5]
\addplot[thick,blue,domain=-7.5:-0.05]{ln(-x)/ln(4)};
\addplot[only marks,mark=*,mark size=1.5pt] coordinates {(-1,0)};
\end{axis}
\end{tikzpicture}
}

\mondai{233}{次の関数の（\quad）内の定義域に対する値域を求めよ．}
\kotae{
対数関数は単調なので，端の値を計算すればよい．底$>1$なら増加，底$<1$なら減少．

(1) $y=\log_{5}x\quad\left(\dfrac{1}{25}\leqq x<5\sqrt{5}\right)$

$x=\dfrac{1}{25}=5^{-2}$：$y=-2$\quad
$x=5\sqrt{5}=5^{3/2}$：$y=\dfrac{3}{2}$（この値は含まない）

$\therefore -2\leqq y<\dfrac{3}{2}$

(2) $y=\log_{1/3}x\quad\left(\dfrac{1}{9}<x<1\right)$

底 $\dfrac{1}{3}<1$ なので単調減少．

$x=\dfrac{1}{9}=\left(\dfrac{1}{3}\right)^{2}$：$y=2$\quad
$x=1$：$y=0$

$\therefore 0<y<2$
}

\mondai{234}{次の数の大小を比べよ．}
\kotae{
同じ底の対数に変換して比較する．底$>1$なら真数が大きいほど対数も大きい．

(1) $\log_{2}3\sqrt{2}=\log_{2}3+\dfrac{1}{2}$，$\log_{2}0.9=\log_{2}\dfrac{9}{10}$，$\log_{2}4=2$

$\log_{2}0.9<0<\log_{2}3\sqrt{2}<\log_{2}4$

$\therefore \log_{2}0.9<\log_{2}3\sqrt{2}<\log_{2}4$

(2) $\log_{1/5}\dfrac{1}{4}$，$\log_{1/5}0.2=\log_{1/5}\dfrac{1}{5}=1$，$\log_{1/5}\dfrac{3}{7}$

底$<1$なので真数が大きいほど対数は小さい．
$\dfrac{1}{5}<\dfrac{1}{4}<\dfrac{3}{7}$ より

$\log_{1/5}\dfrac{3}{7}<\log_{1/5}\dfrac{1}{4}<\log_{1/5}\dfrac{1}{5}=1$
}





\mondai{235}{次の方程式を解け．}
\kotae{
対数方程式は真数を1つにまとめてから指数に戻す．真数条件にも注意．

(1) $\log_{2}(x-3)+\log_{2}(x-1)=3$

真数条件：$x>3$

$\log_{2}(x-3)(x-1)=3$，$(x-3)(x-1)=2^{3}=8$

$x^{2}-4x+3=8$，$x^{2}-4x-5=0$，$(x-5)(x+1)=0$

$x>3$ より $x=5$

(2) $\log_{3}(4x-7)-\log_{3}(x+12)=-1$

真数条件：$x>\dfrac{7}{4}$

$\log_{3}\dfrac{4x-7}{x+12}=-1$，$\dfrac{4x-7}{x+12}=3^{-1}=\dfrac{1}{3}$

$3(4x-7)=x+12$，$12x-21=x+12$，$11x=33$ → $x=3$\quad（$>\dfrac{7}{4}$で適）
}

\mondai{236}{次の不等式を解け．}
\kotae{
底$>1$のとき不等号の向きはそのまま，底$<1$のとき逆転する．
真数条件と連立させる．

(1) $\log_{2}(3x+1)>4$

真数条件：$3x+1>0$，$x>-\dfrac{1}{3}$

底 $2>1$ より $3x+1>2^{4}=16$，$x>5$

真数条件と合わせて $x>5$

(2) $\log_{3}(x-1)<-2$

真数条件：$x>1$

底 $3>1$ より $x-1<3^{-2}=\dfrac{1}{9}$，$x<\dfrac{10}{9}$

$\therefore 1<x<\dfrac{10}{9}$
}

\mondai{237}{$\log_{10}2=0.3010$，$\log_{10}3=0.4771$ とするとき，次の値を求めよ．}
\kotae{
既知の値を組み合わせて表す．

(1) $\log_{10}8=\log_{10}2^{3}=3\times0.3010=0.9030$

(2) $\log_{10}18=\log_{10}(2\cdot3^{2})=\log_{10}2+2\log_{10}3=0.3010+2\times0.4771=1.2552$

(3) $\log_{2}27=\dfrac{\log_{10}27}{\log_{10}2}=\dfrac{3\log_{10}3}{\log_{10}2}=\dfrac{3\times0.4771}{0.3010}=\dfrac{1.4313}{0.3010}\approx4.755$
}

\mondai{238}{次の不等式を満たす最大の整数 $n$ を求めよ．}
\kotae{
(1) $10^{n}\leqq5^{20}$

両辺の常用対数：$n\leqq20\log_{10}5=20(1-\log_{10}2)=20\times0.6990=13.98$

$\therefore n=13$

(2) $10^{-n}\geqq3^{-50}$

$-n\geqq-50\log_{10}3$，$n\leqq50\times0.4771=23.855$

$\therefore n=23$
}

\mondai{239}{ある細菌は1時間で1.5倍に増殖するという．同じ割合で増殖し続けるとき，何時間後にはじめて最初の量の1000倍以上となるか．}
\kotae{
$t$ 時間後の量：$1.5^{t}$ 倍．

$1.5^{t}\geqq1000$ の両辺の常用対数をとる：

$t\log_{10}1.5\geqq3$

$\log_{10}1.5=\log_{10}\dfrac{3}{2}=\log_{10}3-\log_{10}2=0.4771-0.3010=0.1761$

$t\geqq\dfrac{3}{0.1761}\approx17.04$

$\therefore$ 18時間後（はじめて1000倍以上）
}

\end{document}
